% =============================================================
% 第3章 作業の進捗率（WBS・ガントチャート）(課題6) 章本体
% reports.tex から \input される．
% =============================================================

\section{はじめに}
本章には，Arduinoオーケストラ開発における作業進捗を，
WBSタスクごとの進捗率とガントチャートで示す．

\section{WBSタスクごとの進捗率}

\begingroup
\small
\begin{longtable}{p{0.9cm}p{7cm}p{2.5cm}p{2cm}p{2cm}}
\toprule
ID & タスク名 & 担当者 & 進捗率（\%） & 備考 \\
\midrule
\endfirsthead

\toprule
ID & タスク名 & 担当者 & 進捗率（\%） & 備考 \\
\midrule
\endhead

\bottomrule
\endfoot

\multicolumn{5}{l}{\textbf{【100 設計フェーズ ／ 110 共通設計】}} \\
\midrule
111 & FBS作成 & 全員 &  &  \\
112 & PBS作成 & 全員 &  &  \\
113 & WBS作成 & 全員 &  &  \\
114 & MOE・評価基準確定 & 全員 &  &  \\
\midrule
\multicolumn{5}{l}{\textbf{【100 設計フェーズ ／ 120 担当A設計（指揮者人形ユニット）】}} \\
\midrule
121 & サーボ機構設計 & 長見・飛田 & 30\% &  \\
122 & 腕・反射板設計 & 長見・飛田 & 30\% &  \\
123 & LED駆動回路設計 & 長見・飛田 & 30\% &  \\
124 & サーボ駆動回路設計 & 長見・飛田 & 100\% &  \\
125 & サーボ制御SW仕様書 & 長見・飛田 & 100\% &  \\
126 & 輝度エンコードLED制御SW仕様書 & 長見・飛田 & 50\% & 担当Bと要相談 \\
128 & シリアル通信コマンド設計（INSTRUMENT/BPM） & 長見・飛田 & 100\% &  \\
\midrule
\multicolumn{5}{l}{\textbf{【100 設計フェーズ ／ 130 担当B設計（センサー・観測）】}} \\
\midrule
131 & 赤外線センサー配線設計 & 長野・慶田 & 50\% &  \\
132 & フォトTr配線設計 & 長野・慶田 & 80\% &  \\
133 & センサーISR設計 & 長野・慶田 & 100\% &  \\
135 & 状態機械設計 & 長野・慶田 & 100\% &  \\
136 & キャリブレーション手順設計（CALIBRATE状態） & 長野・慶田 & 90\% &  \\
137 & INSTRUMENT\_ID輝度キュー設計 & 長野・慶田 & 80\% &  \\
\midrule
\multicolumn{5}{l}{\textbf{【100 設計フェーズ ／ 140 担当C設計（演奏制御）】}} \\
\midrule
142 & 楽譜進行制御設計 & 長山・慶田 & 80\% &  \\
143 & シリアル送信プロトコル設計 & 長山・慶田 & 90\% &  \\
144 & 楽譜配列定義・5パート & 長山・慶田 & 50\% &  \\
\midrule
\multicolumn{5}{l}{\textbf{【100 設計フェーズ ／ 150 担当D設計（音響合成）】}} \\
\midrule
151 & シリアル受信処理設計 & 中村・尾添 &  &  \\
154 & 音色定義シート作成 & 中村・尾添 &  &  \\
155 & ADSRパラメータ仕様策定 & 中村・尾添 &  &  \\
\midrule
\multicolumn{5}{l}{\textbf{【200 実装フェーズ ／ 210 担当A実装（指揮者人形ユニット）】}} \\
\midrule
211 & サーボ機構組み立て & 長見・飛田 & 0\% &  \\
212 & 腕・反射板取り付け & 長見・飛田 & 0\% &  \\
214 & サーボ駆動回路配線 & 長見・飛田 & 0\% &  \\
215 & サーボ制御SW実装 & 長見・飛田 & 80\% &  \\
216 & 輝度エンコードLED制御SW実装 & 長見・飛田 & 50\% & 担当Bと要相談 \\
217 & Processing GUI実装 & 長見・飛田 & 100\% &  \\
218 & シリアル通信実装（PC→Arduino） & 長見・飛田 & 100\% &  \\
\midrule
\multicolumn{5}{l}{\textbf{【200 実装フェーズ ／ 220 担当B実装（センサー・観測）】}} \\
\midrule
221 & 赤外線センサー配線 & 長野・慶田 & 50\% &  \\
222 & フォトTr配線・遮光チャンバー & 長野・慶田 & 70\% &  \\
223 & センサーISR実装 & 長野・慶田 & 70\% &  \\
224 & テンポ推定実装 & 長野・慶田 & 90\% &  \\
226 & INSTRUMENT\_ID輝度キュー実装 & 長野・慶田 & 50\% &  \\
227 & シリアル通信実装（担当C向け） & 長野・慶田 & 80\% &  \\
\midrule
\multicolumn{5}{l}{\textbf{【200 実装フェーズ ／ 230 担当C実装（演奏制御）】}} \\
\midrule
231 & 拍検出・オフセット実装 & 長山・慶田 & 80\% &  \\
232 & 楽譜進行制御実装 & 長山・慶田 & 80\% &  \\
\midrule
\multicolumn{5}{l}{\textbf{【200 実装フェーズ ／ 240 担当D実装（音響合成）】}} \\
\midrule
241 & シリアル受信処理実装 & 中村・尾添 &  &  \\
242 & 音色合成エンジン実装 & 中村・尾添 &  &  \\
243 & エフェクト・フェルマータ延奏実装 & 中村・尾添 &  &  \\
244 & 音色定義データ作成 & 中村・尾添 &  &  \\
245 & ADSRパラメータ調整 & 中村・尾添 &  &  \\
\midrule
\multicolumn{5}{l}{\textbf{【300 テストフェーズ ／ 310 事前確認テスト（P1-P6）】}} \\
\midrule
311 & 赤外線センサー精度確認 & 長野・慶田 &  &  \\
312 & フォトTr検出距離確認 & 長見・飛田/長野・慶田 &  &  \\
313 & LED輝度・フォトTr検出確認 & 長見・飛田/長野・慶田 &  &  \\
315 & サーボ最大BPM確認 & 長見・飛田 &  &  \\
316 & Arduino-Processing通信確認 & 長野・長山・慶田/中村・尾添 &  &  \\
\midrule
\multicolumn{5}{l}{\textbf{【300 テストフェーズ ／ 320 単体テスト（U1-U4）】}} \\
\midrule
322 & 状態機械単体テスト & 長野・慶田 &  &  \\
323 & 楽譜進行・シリアル送信単体テスト & 長山・慶田 &  &  \\
324 & 音色合成・エフェクト単体テスト & 中村・尾添 &  &  \\
\midrule
\multicolumn{5}{l}{\textbf{【300 テストフェーズ ／ 330 結合テスト（C1-C3）】}} \\
\midrule
331 & Arduino-Processing遅延測定 & 長野・長山・慶田/中村・尾添 &  &  \\
332 & 同期・テンポ追従確認 & 長野・長山・慶田 &  &  \\
333 & フェルマータ動作確認 & 長野・長山・慶田/中村・尾添 &  &  \\
\midrule
\multicolumn{5}{l}{\textbf{【300 テストフェーズ ／ 340 システムテスト（S1-S3）】}} \\
\midrule
341 & 通し演奏テスト & 全員 &  &  \\
342 & 拍タイミング精度評価 & 全員 &  &  \\
343 & テンポ変化・フェルマータ完全確認 & 全員 &  &  \\
\midrule

\end{longtable}
\endgroup

\section{ガントチャート（現在の進捗反映）}

WBSタスクごとの予定週セルを，現時点の進捗率に応じて3段階で色分けして示す．
日付は授業のある水曜日に基づく（6月3日は休講のため6月10日が次の授業日）．
設計フェーズ（100番台）は本ガントチャートには含めない．

\begin{table}[H]
  \centering
  \caption*{凡例（進捗率による色分け）}
  \begin{tabular}{ll@{\hspace{1em}}ll@{\hspace{1em}}ll@{\hspace{1em}}ll}
    \toprule
    \pLow & 80\%未満 & \pMid & 80\%以上100\%未満 & \pDone & 100\% & \pNone & 未着手・未記入 \\
    \bottomrule
  \end{tabular}
\end{table}

\begingroup
\scriptsize
\setlength{\tabcolsep}{3pt}
\begin{longtable}{p{0.7cm}p{3.8cm}p{1.7cm}p{0.8cm}p{0.8cm}p{0.9cm}p{0.8cm}p{0.8cm}p{0.9cm}p{0.9cm}}
\toprule
ID & タスク名 & 担当者 & 進捗 & \makecell{5/27\\(水)} & \makecell{6/10\\(水)} & \makecell{6/17\\(水)} & \makecell{6/24\\(水)} & \makecell{7/1\\(水)} & \makecell{7/8\\(水)} \\
\midrule
\endfirsthead

\toprule
ID & タスク名 & 担当者 & 進捗 & \makecell{5/27\\(水)} & \makecell{6/10\\(水)} & \makecell{6/17\\(水)} & \makecell{6/24\\(水)} & \makecell{7/1\\(水)} & \makecell{7/8\\(水)} \\
\midrule
\endhead

\bottomrule
\endfoot

\multicolumn{10}{l}{\textbf{【200 実装フェーズ ／ 210 担当A実装（指揮者人形ユニット）】}} \\
\midrule
211 & サーボ機構組み立て & 長見・飛田 & 0\% & \pLow & \pLow &  &  &  &  \\
212 & 腕・反射板取り付け & 長見・飛田 & 0\% & \pLow & \pLow &  &  &  &  \\
214 & サーボ駆動回路配線 & 長見・飛田 & 0\% & \pLow & \pLow &  &  &  &  \\
215 & サーボ制御SW実装 & 長見・飛田 & 80\% & \pMid & \pMid &  &  &  &  \\
216 & 輝度エンコードLED制御SW実装 & 長見・飛田 & 50\% & \pLow & \pLow &  &  &  &  \\
217 & Processing GUI実装 & 長見・飛田 & 100\% & \pDone & \pDone &  &  &  &  \\
218 & シリアル通信実装（PC→Arduino） & 長見・飛田 & 100\% & \pDone & \pDone &  &  &  &  \\
\midrule
\multicolumn{10}{l}{\textbf{【200 実装フェーズ ／ 220 担当B実装（センサー・観測）】}} \\
\midrule
221 & 赤外線センサー配線 & 長野・慶田 & 50\% & \pLow & \pLow &  &  &  &  \\
222 & フォトTr配線・遮光チャンバー & 長野・慶田 & 70\% & \pLow & \pLow &  &  &  &  \\
223 & センサーISR実装 & 長野・慶田 & 70\% & \pLow & \pLow &  &  &  &  \\
224 & テンポ推定実装 & 長野・慶田 & 90\% & \pMid &  &  &  &  &  \\
226 & INSTRUMENT\_ID輝度キュー実装 & 長野・慶田 & 50\% & \pLow & \pLow &  &  &  &  \\
227 & シリアル通信実装（担当C向け） & 長野・慶田 & 80\% & \pMid & \pMid &  &  &  &  \\
\midrule
\multicolumn{10}{l}{\textbf{【200 実装フェーズ ／ 230 担当C実装（演奏制御）】}} \\
\midrule
231 & 拍検出・オフセット実装 & 長山・慶田 & 80\% & \pMid & \pMid &  &  &  &  \\
232 & 楽譜進行制御実装 & 長山・慶田 & 80\% & \pMid & \pMid &  &  &  &  \\
\midrule
\multicolumn{10}{l}{\textbf{【200 実装フェーズ ／ 240 担当D実装（音響合成）】}} \\
\midrule
241 & シリアル受信処理実装 & 中村・尾添 & — & \pNone & \pNone &  &  &  &  \\
242 & 音色合成エンジン実装 & 中村・尾添 & — & \pNone & \pNone &  &  &  &  \\
243 & エフェクト・フェルマータ延奏実装 & 中村・尾添 & — & \pNone & \pNone &  &  &  &  \\
244 & 音色定義データ作成 & 中村・尾添 & — & \pNone & \pNone &  &  &  &  \\
245 & ADSRパラメータ調整 & 中村・尾添 & — & \pNone & \pNone &  &  &  &  \\
\midrule
\multicolumn{10}{l}{\textbf{【300 テストフェーズ ／ 310 事前確認テスト（P1-P6）】}} \\
\midrule
311 & 赤外線センサー精度確認 & 長野・慶田 & — & \pNone & \pNone &  &  &  &  \\
312 & フォトTr検出距離確認 & 長見・飛田/長野・慶田 & — & \pNone & \pNone &  &  &  &  \\
313 & LED輝度・フォトTr検出確認 & 長見・飛田/長野・慶田 & — & \pNone & \pNone &  &  &  &  \\
315 & サーボ最大BPM確認 & 長見・飛田 & — & \pNone & \pNone &  &  &  &  \\
316 & Arduino-Processing通信確認 & 長野・長山・慶田/中村・尾添 & — & \pNone & \pNone &  &  &  &  \\
\midrule
\multicolumn{10}{l}{\textbf{【300 テストフェーズ ／ 320 単体テスト（U1-U4）】}} \\
\midrule
322 & 状態機械単体テスト & 長野・慶田 & — & \pNone &  & \pNone &  &  &  \\
323 & 楽譜進行・シリアル送信単体テスト & 長山・慶田 & — & \pNone &  & \pNone &  &  &  \\
324 & 音色合成・エフェクト単体テスト & 中村・尾添 & — & \pNone &  & \pNone &  &  &  \\
\midrule
\multicolumn{10}{l}{\textbf{【300 テストフェーズ ／ 330 結合テスト（C1-C3）】}} \\
\midrule
331 & Arduino-Processing遅延測定 & 長野・長山・慶田/中村・尾添 & — & \pNone &  & \pNone &  &  &  \\
332 & 同期・テンポ追従確認 & 長野・長山・慶田 & — & \pNone &  & \pNone & \pNone &  &  \\
333 & フェルマータ動作確認 & 長野・長山・慶田/中村・尾添 & — & \pNone &  & \pNone & \pNone &  &  \\
\midrule
\multicolumn{10}{l}{\textbf{【300 テストフェーズ ／ 340 システムテスト（S1-S3）】}} \\
\midrule
341 & 通し演奏テスト & 全員 & — & \pNone &  &  & \pNone &  &  \\
342 & 拍タイミング精度評価 & 全員 & — & \pNone &  &  & \pNone &  &  \\
343 & テンポ変化・フェルマータ完全確認 & 全員 & — & \pNone &  &  & \pNone &  &  \\
\midrule

\end{longtable}
\endgroup
